\begin{frame}[t]{\ev{\tagM\,\tagB\,\tagP}Separate calls represent separate decisions}
\vspace{0pt}
{\small
\begin{tabularx}{\textwidth}{@{}>{\ttfamily}L{2.2cm} L{5.4cm} Y@{}}
\toprule
\normalfont Call & Responsibility & Status\\
\midrule
checkpoint & identify the captured state & \tagM\ named snapshots exercised\\
fork & create a child with declared sharing & \tagP\\
evaluate & produce artifact-bound evidence & \tagP\ external evaluator\\
select & choose one candidate & \tagP\\
reduce & combine measurements & \tagB\ deduplication\\
promote & authorise publication & \tagB\ digest gate\\
\bottomrule
\end{tabularx}\par}
\vspace{.1cm}
{\small
\begin{ticks}
\item fork declares the seed policy, the lease identity and the generation
\item[\no] the built reducer carries lineage, without dependence-adjusted uncertainty or abstention
\item a fresh test is a separate evaluation event, outside select
\end{ticks}\par}
\claim{Select chooses a candidate. Reduce combines measurements.}
\sources{Reducer: Eliaz, \href{https://arxiv.org/abs/2607.09689}{arXiv:2607.09689}. Fork contract and interface: our design documents.}
\note{[0:40] We separate six calls, because each one is a different decision. Checkpoint identifies captured state, and our named snapshots exercised it. Fork, evaluate and select are proposed. Fork must declare its seed policy, lease identity and generation. Reduce has built deduplication. It carries lineage but does not yet adjust uncertainty for dependence or abstain. Promote has a built digest gate. Next, we place these calls in one architecture.}
\end{frame}
